\section{Least viable dimensions and a finite dyadic index}

The low-density argument of~\cite[Proposition~prop:albootstrap-ld,
lines 1112--1128]{Arad2016RigorousRG} uses the least dimension of a
viable subspace for each spectral target. The following results concern
these natural minima and the scalar count used in the argument.

\begin{theorem}[Restriction of a viable target]
    \label{thm:viable_target_restriction}
    \lean{Submodule.IsViable.of_target_le}
    \leanok
    Let $M,F$ be finite-dimensional complex inner-product spaces.
    If $S\subseteq M$ is $\delta$-viable for $T_+$ and
    $T_-\subseteq T_+\subseteq M\otimes F$, then $S$ is
    $\delta$-viable for $T_-$.
\end{theorem}
\begin{proof}
    \leanok
    \uses{def:subspace_viable,thm:subspace_close_iff}
    The inequality
    $\|P_{S\otimes F}x\|^2\ge(1-\delta)\|x\|^2$ for $x\in T_+$
    remains valid after restricting $x$ to $T_-$.
\end{proof}

\begin{theorem}[An attained least viable dimension]
    \label{thm:viable_minimum}
    \lean{Submodule.exists_isViable_minimal_finrank}
    \leanok
    Let $\delta\ge0$ and $T\subseteq M\otimes F$.
    There are a natural number $p\le\dim M$ and a subspace
    $W\subseteq M$ of dimension $p$ that is $\delta$-viable for $T$,
    such that $p\le\dim S$ for every $\delta$-viable subspace $S$ for $T$.
\end{theorem}
\begin{proof}
    \leanok
    The set of dimensions of viable subspaces is a nonempty subset of
    the natural numbers: it contains $\dim M$, since $M$ is viable.
    Its least member is the dimension of an actual viable subspace.
\end{proof}

\begin{theorem}[Increasing minima for increasing targets]
    \label{thm:viable_minimum_sequence}
    \lean{Submodule.exists_isViable_minimal_sequence}
    \leanok
    Let $\delta\ge0$ and $T_j\subseteq T_{j+1}\subseteq M\otimes F$.
    One can choose, for every $j\ge0$, a $\delta$-viable subspace
    $W_j$ for $T_j$ whose dimension $p_j$ is least among all such
    subspaces. These choices satisfy $p_j\le\dim M$ and
    $p_j\le p_{j+1}$.
\end{theorem}
\begin{proof}
    \leanok
    \uses{thm:viable_minimum,thm:viable_target_restriction}
    Choose an attained minimum for each target. The space $W_{j+1}$
    is also viable for $T_j$, so the minimum property gives
    $p_j\le\dim W_{j+1}=p_{j+1}$.
\end{proof}

\begin{theorem}[A finite index exceeding a prescribed dimension]
    \label{thm:dyadic_filter_count_finite_index}
    \lean{Real.exists_dyadic_polynomial_filter_count_ge}
    \leanok
    Fix $C\ge1$, $v>0$ and an integer $m\ge2$. For $j\ge0$, put
    \begin{align*}
      g_j&=\min\{v/(6\cdot2^j),1\},&
      u_j&=1+\log m+\log(1/g_j),\\
      \ell_j&=\lceil Cu_j^2/g_j\rceil,&
      D_j&=(m\ell_j^2)^{22\ell_j}.
    \end{align*}
    For every natural number $N$ there is a finite index $j$ such that
    $N\le\ell_j$ and $N\le D_j$.
\end{theorem}
\begin{proof}
    \leanok
    \uses{lem:ceil_quadratic_scale}
    Take $j=\lceil Nv/(6C)\rceil$. Since $j\le2^j$, one has
    $Nv/(6\cdot2^j)\le C$. Hence
    $N\le C/g_j\le Cu_j^2/g_j\le\ell_j$.
    Finally $\ell_j\ge1$ and
    $\ell_j\le\ell_j^2\le m\ell_j^2\le D_j$.
\end{proof}
